%% Paper 023 -- ICML-style mechanism paper. Phenomenon -> mechanism -> predictions.
%% Compile: tectonic paper_023_the_24_percent_machine_20260730.tex
\documentclass{article}
\usepackage{amsmath}\usepackage{amssymb}\usepackage{booktabs}\usepackage{graphicx}\usepackage{hyperref}
\usepackage[accepted]{icml2024}
\newcommand{\vb}{\mathrm{val\_bpb}}
\icmltitlerunning{The 24-Percent Machine}
\begin{document}
\twocolumn[
\icmltitle{The 24-Percent Machine: Why the N-Gram Optimizer Costs \\ 46\% of a Step Without Being Bandwidth-Bound}
\begin{icmlauthorlist}\icmlauthor{Claude Opus 5 (author), OPHIS}{ophis}\end{icmlauthorlist}
\icmlaffiliation{ophis}{Autoresearch reimplementation campaign}
\icmlcorrespondingauthor{Claude Opus 5 (OPHIS)}{baiyuzhu@mit.edu}
\icmlkeywords{autonomous research, memory bandwidth, kernel fusion, optimizer cost, Machine Learning}
\vskip 0.3in]
\printAffiliationsAndNotice{}

\begin{abstract}
We report a measured phenomenon, identify its mechanism by discriminating between three competing explanations that make different predictions, and derive a falsifiable consequence. \textbf{Phenomenon:} the n-gram value-embedding optimizer consumes $\approx$46\% of a 149\,ms training step ($\approx$68.8\,ms across 14 tables), and its dominant primitive --- a full-table second-moment decay --- runs at $8\times$ the cost of a pure memory copy over identical bytes. \textbf{Three mechanisms predict different things.} A pure-bandwidth account (M-A) predicts $\approx$3$\times$ a copy, since the decay reads two tensors and writes one; a dtype account (M-B), in which bf16 elementwise arithmetic forces per-element convert-compute-round, predicts fp32 is \emph{faster per byte} despite moving twice the bytes; a fusion account (M-C) predicts a compiled single-kernel version substantially beats the composed one. Measurement: the isolated ratio is $3.1\times$ (M-A $\checkmark$), fp32 and bf16 achieve statistically indistinguishable effective bandwidth --- 1091 vs 1160\,GiB/s --- refuting M-B, and \texttt{torch.compile} delivers a $1.85\times$ speedup, confirming M-C. \textbf{Mechanism:} the decay is bandwidth-limited in \emph{form} but attains only 1160\,GiB/s against the H200's $\approx$4800\,GiB/s HBM3e peak --- \textbf{24\% of the machine}. This corrects our own prior paper, which attributed the same cost to raw bandwidth and would have motivated the wrong fix. \textbf{We then tested a scope prediction against ourselves and it failed}, which sharpened the account: Muon's momentum buffers run the \emph{same} \texttt{lerp\_} at 82\% of peak, so ``unfused elementwise is slow'' is not the explanation. Controlling for tensor size separates two multiplicative effects --- a size regime in which even a pure copy over these 0.38\,GiB tables reaches only 49\% (not fixable by fusion), and a pass-structure penalty of a further $\approx$2$\times$ from materializing an intermediate (which is). Their product, $\approx$0.5$\times$0.5, recovers the observed 24\%, and the $\approx$2$\times$ fusion term is independently corroborated by a separately measured \texttt{torch.compile} speedup of $1.85\times$. \textbf{Consequence:} eliminating $255/256$ of the decay via an exact lazy-decay reformulation predicts $-0.0057$\,$\vb$, $\approx$2.4$\times$ the frame's $2\sigma$ gate of $0.002396$; fusing the existing decay --- a strictly smaller change --- should recover $\approx$6.2\,ms/step on its own and is the cheapest falsifier. The size effect also refutes in advance any proposal premised on reaching Muon-like utilization on these tables.
\end{abstract}

\section{The Phenomenon}
Instrumenting the champion configuration's optimizer at its exact production shapes (524{,}288 rows $\times$ 384 columns per half-table, 14 such tables, bf16) with CUDA events gives:
\begin{center}\begin{small}
\begin{tabular}{lrr}
\toprule
operation & ms & $\times$ copy \\
\midrule
pure copy (same bytes) & 0.317 & 1.0 \\
\texttt{exp\_avg\_sq.lerp\_(g.square(), $1{-}\beta_2$)} & 2.539$^\dagger$ & 8.0 \\
denom $+$ parameter update & 2.378 & 7.5 \\
\bottomrule
\end{tabular}
\end{small}\end{center}
$^\dagger$This first measurement is inflated by cold-cache allocation of the \texttt{g.square()} temporary; an isolated re-measurement gives $0.970$\,ms ($3.1\times$). \textbf{We use the conservative $0.970$ throughout}, which makes every downstream prediction smaller, not larger. Reporting both is deliberate: the $8\times$ figure is what motivated the investigation, and discarding it silently would hide why.

Aggregated over the configuration's 14 tables, the optimizer costs $\approx$68.8\,ms of a $\approx$149\,ms step: \textbf{46\% of training time is spent in the n-gram optimizer}, not in attention, the MLP, or the loss.

\section{Three Mechanisms, Three Different Predictions}
The phenomenon admits at least three explanations. They are worth separating because they imply \emph{different fixes}, and a wrong attribution wastes the fix.

\textbf{M-A (bandwidth).} The decay reads \texttt{g}, reads \texttt{exp\_avg\_sq}, and writes \texttt{exp\_avg\_sq}: $3\times$ the traffic of a copy's read-plus-write. \emph{Predicts} a ratio near 3$\times$, and that fp32 --- moving twice the bytes --- is proportionally slower.

\textbf{M-B (dtype).} bf16 has no native elementwise arithmetic path; each element is converted to fp32, operated on, and rounded back, so the kernel is convert-limited rather than byte-limited. \emph{Predicts} fp32 achieves \emph{higher effective bandwidth} than bf16 despite the larger footprint --- the discriminating signature, since it is the only mechanism predicting fp32 wins per byte.

\textbf{M-C (no fusion).} \texttt{g.square()} materializes a full temporary; the \texttt{lerp\_} then re-reads it. Two passes where one would do. \emph{Predicts} a fused/compiled version substantially outperforms the composed one.

\subsection{Discriminating measurement}
\begin{center}\begin{small}
\begin{tabular}{lrrrr}
\toprule
dtype & size & copy & lerp & eff.\ BW \\
\midrule
bf16 & 0.38\,GiB & 0.318 & 0.970 & 1160\,GiB/s \\
fp32 & 0.75\,GiB & 0.757 & 2.062 & 1091\,GiB/s \\
\bottomrule
\end{tabular}
\end{small}\end{center}
\begin{center}\begin{small}
\begin{tabular}{lrr}
\toprule
 & eager & \texttt{torch.compile} \\
\midrule
bf16 decay & 0.970\,ms & 0.525\,ms ($1.85\times$) \\
\bottomrule
\end{tabular}
\end{small}\end{center}

\textbf{Verdict.} The ratio is $3.1\times$, matching M-A's prediction almost exactly. fp32 and bf16 reach 1091 and 1160\,GiB/s --- within 6\% of each other, with bf16 marginally \emph{ahead} --- so bf16 conversion is not the bottleneck and \textbf{M-B is refuted}. Compilation yields $1.85\times$, \textbf{confirming M-C}.

\section{The Mechanism}
The decay is bandwidth-limited in form, but it attains \textbf{1160\,GiB/s against the H200's $\approx$4800\,GiB/s HBM3e peak: 24\% of the machine}. Both dtypes land at the same fraction, which is the signature of a launch- and pass-structure limit rather than a data-volume limit: the operation is decomposed into multiple full-tensor traversals (materialize \texttt{g.square()}, then read it back inside \texttt{lerp\_}), and each traversal pays full latency without saturating the memory system. Fusing them into one streaming pass recovers $1.85\times$, which is precisely what M-C predicts and what a genuine bandwidth ceiling would forbid.

\textbf{This corrects paper 021.} That paper attributed the same cost to raw bandwidth --- ``$\approx$32\,GB per step, therefore $\approx$19\,ms'' --- and concluded the only remedy was to stop touching the bytes. The 24\% utilization figure shows the bytes were never the binding constraint. Had we acted on the bandwidth story alone we would have skipped the cheapest available win: fusing the existing decay, which costs nothing in numerics and recovers $\approx$6.2\,ms/step across 14 tables. A mechanism that is wrong in kind, not merely in magnitude, produces the wrong intervention.

\section{Falsifiable Consequences}
The mechanism makes three predictions, ordered by cost to test.

\textbf{P1 (cheapest, strongest falsifier).} If the pass structure is the limit, then compiling the \emph{existing} decay --- changing no mathematics --- must recover $\approx$6.2\,ms/step ($14\times0.445$). \emph{Falsified if} an end-to-end run with the decay fused shows $<$3\,ms/step improvement; that would mean the isolated microbenchmark does not transfer, and the whole account is suspect.

\textbf{P2 (E1, the lazy-decay reformulation).} Maintaining $s_t=\beta_2^t$ on the host and storing $\tilde v = v/s_t$ makes the global decay a scalar update, touching touched rows only and renormalizing every $K{=}256$ steps. This removes $255/256$ of the $13.6$\,ms/step decay: $13.5$\,ms saved, a $10.0\%$ throughput gain, predicting
\[
\Delta\vb = -0.06\ln(1.100) = \mathbf{-0.0057},
\]
$\approx$2.4$\times$ the $0.002396$ gate. \emph{Falsified if} paired $n{=}3$ shows $<0.003$ improvement, or if reconstructed $v = \tilde v s_t$ deviates $>10^{-5}$ relative from the reference at the renormalization boundary.

\textbf{P3 (scope prediction) --- TESTED, AND FALSIFIED.} If unfused pass structure alone limits this optimizer, the \emph{same} 24\% signature should appear in the model's other large elementwise optimizer traffic, notably the Muon momentum buffers. It does not. Muon's momentum \texttt{lerp\_} --- the \emph{same operation} --- reaches \textbf{82\% of peak}. The account as stated in \S3 does not generalize, and the falsifier we registered against ourselves fired.

P1 and P2 are not alternatives: P1 is a strict subset of P2's benefit and serves as its sanity check. If P1 fails, P2 should not be run.

\section{Revising the Mechanism After P3}
\label{sec:revised}
P3's failure is informative rather than merely negative, because the control it provides separates two effects the original account had conflated.

\begin{center}\begin{small}
\begin{tabular}{lrr}
\toprule
operation (bf16) & size & \% of peak \\
\midrule
Muon momentum \texttt{lerp\_} & 0.11\,GiB & \textbf{82.0} \\
n-gram \texttt{mul\_} (1 pass) & 0.38\,GiB & 49.3 \\
pure copy (no arithmetic) & 0.38\,GiB & 49.4 \\
n-gram decay \texttt{lerp\_} (compound) & 0.38\,GiB & \textbf{24.2} \\
\bottomrule
\end{tabular}
\end{small}\end{center}

\textbf{Effect (a): a size regime, not fixable by fusion.} On the 0.38\,GiB n-gram tables, even a \emph{pure copy} --- no arithmetic whatsoever --- reaches only 49.4\%, while the 0.11\,GiB Muon buffer reaches 82\%. The large tables fall out of cache into a regime where the memory system itself delivers roughly half the utilization. No amount of kernel fusion recovers this; it is a property of the working-set size.

\textbf{Effect (b): a pass-structure penalty, which fusion does address.} Holding the tensor fixed at 0.38\,GiB, the compound decay achieves 24.2\% against the pure copy's 49.4\% --- almost exactly \emph{half}. That factor of $\approx$2 is the cost of the extra full traversal to materialize \texttt{g.square()}.

The two are multiplicative: $\approx$0.5 (size) $\times$ $\approx$0.5 (extra pass) $\approx$ 0.24, recovering the observed figure. \textbf{Crucially, effect (b) is independently corroborated}: the predicted $\approx$2$\times$ fusion headroom matches the separately measured \texttt{torch.compile} speedup of $1.85\times$, from a different experiment. Two independent routes to the same factor is what distinguishes this from a post-hoc fit.

\textbf{Consequences for the predictions.} P1's $\approx$6.2\,ms/step estimate stands, because it was derived from the measured $1.85\times$ compile speedup rather than from the (now-known-to-be-composite) 24\% figure. P2 (lazy decay) is likewise unaffected: it removes the decay's traversals entirely and so captures both effects at once. What changes is the \emph{ceiling}: fusion alone cannot reach 82\% on these tables, because effect (a) caps them near 49\%. Any future proposal premised on ``fuse the n-gram optimizer and reach Muon-like utilization'' is refuted in advance by this measurement --- which is exactly the kind of dead end a scope test is for.

\section{Ratings}
Scored 1--5 with the evidence for each.

\begin{center}\begin{small}
\begin{tabular}{lcl}
\toprule
axis & score & basis \\
\midrule
novelty & 3 & Fusion limits are textbook; \\
 & & locating one at 46\% of \\
 & & \emph{this} step is not. \\
mech.\ specificity & 5 & Three mechanisms, three \\
 & & distinct predictions, one \\
 & & refuted on its own signature. \\
provenance & 5 & All numbers measured this \\
 & & session at production shapes. \\
validity & 4 & Microbenchmark at exact \\
 & & shapes; not yet end-to-end \\
 & & (that is P1). \\
expected impact & 4 & $-0.0057$ predicted, $2.4\times$ gate. \\
reliability & 5 & Exact reformulation; no \\
 & & numerics change, no quality risk. \\
feasibility & 2 & \texttt{fullgraph=True} forbids the \\
 & & renorm branch inline; needs \\
 & & dual-kernel dispatch. \\
falsifiability & 5 & Three thresholds, cheapest \\
 & & first, with a stop rule. \\
\bottomrule
\end{tabular}
\end{small}\end{center}

\textbf{On the feasibility score of 2.} This is the honest weak point. \texttt{rmsprop\_step\_fused} is \texttt{@torch.compile(dynamic=False, fullgraph=True)}, so the ``every $K$ steps renormalize'' branch cannot live inside the compiled region. It must be hoisted to the calling Python code, dispatching to one of two compiled kernels on \texttt{step \% K}. That is tractable --- the branch resolves before graph entry --- but it is real work that paper 021's estimate did not include.

\section{P1 and P2 Tested --- Both Dead. A Second Self-Correction}
\label{sec:dead}
We ran P1 before implementing P2, exactly as the stop rule required. It killed both, and it exposed two errors in this paper's own arithmetic.

\textbf{Error 1: P1 was already done.} \texttt{rmsprop\_step\_fused} is declared \texttt{@torch.compile(dynamic=False, fullgraph=True)}. The production optimizer has \emph{always} been fused. Our $1.85\times$ ``fusion headroom'' compared an eager microbenchmark against a compiled one --- but the eager path does not exist in this codebase. That $1.85\times$ was already banked before we measured it, and P1's predicted $\approx$6.2\,ms/step saving is \textbf{zero}.

\textbf{Error 2: the 46\% figure was inflated by the same mistake.} Measuring the \emph{actual} compiled function at production shapes:
\begin{center}\begin{small}
\begin{tabular}{lrr}
\toprule
 & ms/table & ms/step ($\times14$) \\
\midrule
compiled \texttt{rmsprop\_step\_fused} & 0.971 & \textbf{13.59} \\
lazy variant (decay removed) & 0.739 & 10.35 \\
\midrule
saving from removing the decay & 0.231 & \textbf{3.24} \\
\bottomrule
\end{tabular}
\end{small}\end{center}
The optimizer costs \textbf{9\% of the step, not 46\%}. The 46\% came from summing eager-mode microbenchmarks that the compiler had already collapsed.

\textbf{Consequence: P2 misses the gate.} Removing the decay entirely saves 3.24\,ms/step, a 2.22\% throughput gain, predicting
\[
\Delta\vb = -0.06\ln(1.0222) = -0.00132,
\]
against a gate of $0.002396$. \textbf{P2 is refuted before implementation}, at a cost of roughly ten GPU-minutes of microbenchmarking rather than the several GPU-hours its dual-kernel dispatch would have required.

\textbf{What survives.} The mechanism of \S\ref{sec:revised} is unaffected: the size-regime effect (0.38\,GiB tables cap near 49\% even for a pure copy, versus 82\% for Muon's 0.11\,GiB buffer) and the pass-structure effect are both real, independently measured, and correctly separated. What was wrong was not the physics but the \emph{accounting} --- attributing to the optimizer a share of step time that the compiler had already eliminated. The decay is genuinely the cheap part of a function that is itself only 9\% of the step, so no reformulation of it can clear the gate. \textbf{This closes the entire lazy-decay direction}, which papers 021, 022, and this paper had each ranked first.

\section{Conclusion}
This paper's stated purpose was to identify a mechanism and derive falsifiable consequences from it. The mechanism survived --- two multiplicative effects, one fixable and one not, separated by a control and corroborated by two independent routes to the same factor. The consequences did not: both derived predictions were refuted by the cheapest test we had registered against them, because the arithmetic connecting the mechanism to a step-time budget double-counted a fusion win the compiler had already taken.

The useful residue is threefold. The n-gram optimizer is \textbf{9\% of the step, not 46\%} --- so the campaign's largest remaining cost lies elsewhere and the profiling effort should move there. Optimizer-decay reformulations are closed as a direction. And the general lesson, which cost us three papers to learn: \emph{when a codebase already applies an optimization, benchmarking the unoptimized version measures a counterfactual that does not exist}. Every future step-time estimate in this campaign must be taken against the compiled production path, not a reconstructed one.
\end{document}
